\documentclass{article}
\usepackage{../assignment_preamble}

\title{Homework 2}
\author{Ravi Kini}
\date{October 17, 2025}

\begin{document}

\maketitle

\problem{}
\begin{align*}
    \begin{bmatrix}
        1 & 2 & 7 \\
        3 & 5 & -1 \\
        6 & 1 & 4
    \end{bmatrix}
\end{align*}
Putting the matrix in row echelon form using Gaussian elimination:
\begin{align*}
    & \begin{gmatrix}[b]
        1 & 2 & 7 \\
        3 & 5 & -1 \\
        6 & 1 & 4
    \rowops
    \add[-3][]{0}{1}
    \add[-6][]{0}{2}
    \end{gmatrix}  \\
    & \begin{gmatrix}[b]
        1 & 2 & 7 \\
        0 & -1 & -22 \\
        0 & -11 & -38
    \rowops
    \add[-11][]{1}{2}
    \end{gmatrix}   \\
    & \begin{gmatrix}[b]
        1 & 2 & 7 \\
        0 & -1 & -22 \\
        0 & 0 & 204
    \end{gmatrix} 
\end{align*}
The LU factorization of the matrix is then:
\begin{align*}
    \begin{bmatrix}
        1 & 2 & 7 \\
        3 & 5 & -1 \\
        6 & 1 & 4
    \end{bmatrix}
    & =
    \begin{bmatrix}
        1 & 0 & 0 \\
        3 & 1 & 0 \\
        6 & 11 & 1
    \end{bmatrix}
    \begin{bmatrix}
        1 & 2 & 7 \\
        0 & -1 & -22 \\
        0 & 0 & 204
    \end{bmatrix}
\end{align*}

\clearpage

\problem{}
Let $A \in \mathbb{C}^{n\times n}$ and $b, x \in \mathbb{C}^n$. Solving for $x$ in $Ax = b$ can be done using only real-valued matrices and vectors as follows:

Note that $A$ can be decomposed into two real-valued matrices as $A_{\textrm{Re}} + iA_{\textrm{Im}}$. Similarly, $b$ and $x$ can be decomposed as $b_{\textrm{Re}} + ib_{\textrm{Im}}$ and $x_{\textrm{Re}} + ix_{\textrm{Im}}$. Then $Ax = b$ can be rewritten as:
\begin{align*}
    (A_{\textrm{Re}} + iA_{\textrm{Im}})(x_{\textrm{Re}} + ix_{\textrm{Im}}) & = b_{\textrm{Re}} + ib_{\textrm{Im}} \\
    (A_{\textrm{Re}}x_{\textrm{Re}} - A_{\textrm{Im}}x_{\textrm{Im}}) + i(A_{\textrm{Re}}x_{\textrm{Im}} + A_{\textrm{Im}}x_{\textrm{Re}}) & = 
\end{align*}
This can be modelled as the following linear system, which uses only real-valued matrices:
\begin{align*}
    \begin{bmatrix}
        A_{\textrm{Re}} & -A_{\textrm{Im}} \\
        A_{\textrm{Im}} & A_{\textrm{Re}}
    \end{bmatrix}
    \begin{bmatrix}
        x_{\textrm{Re}} \\ x_{\textrm{Im}}
    \end{bmatrix}
    & = 
    \begin{bmatrix}
        b_{\textrm{Re}} \\ b_{\textrm{Im}}
    \end{bmatrix}
\end{align*}

\clearpage

\problem{}
Let $L$ be an invertible $2\times 2$ lower triangular matrix. It then has the form:
\begin{align*}
    L & = \begin{bmatrix}
        l_{1,1} & 0 \\
        l_{2,1} & l_{2,2} \\
    \end{bmatrix}
\end{align*}
Its inverse $L^{-1}$ then has the form:
\begin{align*}
    L^{-1} & = (l_{1,1}l_{2,2})^{-1}\begin{bmatrix}
        l_{2,2} & 0 \\
        -l_{2,1} & l_{1,1} \\
    \end{bmatrix}
\end{align*}
$L^{-1}$ is then also a lower triangular matrix. Now suppose that for all invertible $n\times n$ lower triangular matrices, the inverse is also lower triangular. Then let $L$ be an invertible $(n + 1) \times (n + 1)$ lower triangular matrix. It then has the form:
\begin{align*}
    L & = \begin{bmatrix}
        L' & 0 \\
        L'' & l_{n + 1, n + 1} \\
    \end{bmatrix}
\end{align*}
where $L'$ is an invertible $n \times n$ lower triangular matrix and $L''$ is a $1\times n$ matrix. Its inverse $L^{-1}$ then has the form:
\begin{align*}
    L^{-1} & = \begin{bmatrix}
        (L')^{-1} & 0 \\
        -l_{n + 1, n + 1}^{-1}L''{L'}^{-1} & l_{n + 1, n + 1}^{-1} \\
    \end{bmatrix}
\end{align*}
$L^{-1}$ is then also a lower triangular matrix. By induction, for all invertible lower triangula matrices, the inverse is lower triangular.

\clearpage

\problem{}
\begin{align*}
    M & = \begin{bmatrix}
        A & B \\ C & D
    \end{bmatrix}
\end{align*}
\subproblem{(a)}
Consider the matrix product:
\begin{align*}
    \begin{bmatrix}
        I & 0 \\ CA^{-1} & I
    \end{bmatrix}
    \begin{bmatrix}
        A & 0 \\ 0 & D - CA^{-1}B
    \end{bmatrix}
    \begin{bmatrix}
        I & A^{-1}B \\ 0 & I
    \end{bmatrix}
\end{align*}
It can be simplified as follows:
\begin{align*}
    \begin{bmatrix}
        I & 0 \\ CA^{-1} & I
    \end{bmatrix}
    \begin{bmatrix}
        A & 0 \\ 0 & D - CA^{-1}B
    \end{bmatrix}
    \begin{bmatrix}
        I & A^{-1}B \\ 0 & I
    \end{bmatrix} & = \begin{bmatrix}
        A & 0 \\ C & D - CA^{-1}B
    \end{bmatrix}
    \begin{bmatrix}
        I & A^{-1}B \\ 0 & I
    \end{bmatrix} \\
    & = \begin{bmatrix}
        A & B \\ C & D
    \end{bmatrix} = M
\end{align*}

\subproblem{(b)}
Let $A$ have the LU factorization $L_1U_1$ and $D - CA^{-1}B$ have the LU factorization $L_2U_2$. $M$ can then be rewritten as:
\begin{align*}
    M & = \begin{bmatrix}
        I & 0 \\ CA^{-1} & I
    \end{bmatrix}
    \begin{bmatrix}
        L_1U_1 & 0 \\ 0 & L_2U_2
    \end{bmatrix}
    \begin{bmatrix}
        I & A^{-1}B \\ 0 & I
    \end{bmatrix} \\
    & = \begin{bmatrix}
        I & 0 \\ CA^{-1} & I
    \end{bmatrix}
    \begin{bmatrix}
        L_1 & 0 \\ 0 & L_2
    \end{bmatrix}
    \begin{bmatrix}
        U_1 & 0 \\ 0 & U_2
    \end{bmatrix}
    \begin{bmatrix}
        I & A^{-1}B \\ 0 & I
    \end{bmatrix} \\
    & = \begin{bmatrix}
        L_1 & 0 \\ CA^{-1}L_1 & L_2
    \end{bmatrix}
    \begin{bmatrix}
        U_1 & U_1A^{-1}B \\ 0 & U_2
    \end{bmatrix} \\
\end{align*}
The matrices are lower triangular and upper triangular, respectively. The LU factorization of $M$ is then:
\begin{align*}
    M & = \begin{bmatrix}
        L_1 & 0 \\ CA^{-1}L_1 & L_2
    \end{bmatrix}
    \begin{bmatrix}
        U_1 & U_1A^{-1}B \\ 0 & U_2
    \end{bmatrix} \\
\end{align*}
\subproblem{(c)}
The following recursive algorithm finds the LU factorization of a matrix $M$:\
\begin{algorithm}
\caption{LU}
\begin{algorithmic}
\State $n \gets \textrm{size}(M)$
\State $\begin{bmatrix} A & B \\ C & D \end{bmatrix} \gets M$
\State $L_1, U_1 \gets \textrm{LU}(A)$
\State $L_2, U_2 \gets \textrm{LU}(D - CA^{-1}B)$
\State $L \gets \begin{bmatrix} L_1 & 0 \\ CA^{-1}L_1 & L_2 \end{bmatrix}$
\State $U \gets \begin{bmatrix}U_1 & U_1A^{-1}B \\ 0 & U_2\end{bmatrix}$
\end{algorithmic}
\end{algorithm}
LU factorization vis Gaussian elimination has a time complexity of $O(n^3)$. Block LU factorization also has a time complexity of $O(n^3)$, mainly due to having to perform matrix inversion.

\clearpage

\problem{}
Suppose $A$ has two different LU factorizations, $L_1U_1$ and $L_2U_2$. Then:
\begin{align*}
    L_1U_1 & = L_2U_2 \\
    L_2^{-1}L_1U_1 & = U_2 \\
    L_2^{-1}L_1 & = U_2U_1^{-1} \\
\end{align*}
Since the product of two unit lower triangular matrices is unit lower triangular, and likewise for upper triangular matrices, $L_2^{-1}L_1 = U_2U_1^{-1}$ must be both upper and lower triangular. The only matrix that is both upper and lower triangular is $I$, which means that $L_2^{-1}L_1 = U_2U_1^{-1} = I$. For this to be true, $L_1 = L_2$ and $U_1 = U_2$; consequently, each matrix permits a unique LU factorization.

\clearpage

\problem{}
\subproblem{(a)}
Suppose $A$ is symmetric positive definite. Let $A$ have the LU factorization $LU$. Since $A$ is symmetric $A = A^T$, and so $A = LU = U^TL^T$. $U^T$ is lower triangular; for some diagonal $D$, it can be written as ${U'}^TD$, where ${U'}^T$ is unit lower triangular. ${L'}^T = DL^T$ is then upper triangular, so ${U'}^T{L'}^T$ is also a LU factorization of $A$. By the uniqueness of the LU factorization, $L = {U'}^T$ and $U = {L'}^T$. Since $A$ is positive definite, for any $x \neq 0$:
\begin{align*}
    x^TAx & = x^T(LDL^T)x \\
    & = (L^Tx)^T D (L^Tx) \\
    & = y^T Dy > 0
\end{align*} 
For this to be the case, the diagonal entries of $D$ must be positive. We can then define $D'$ such that ${D'}^2 = D$. Then:
\begin{align*}
    A & = LDL^T \\
    & = L{D'}^2L^T \\
    & = (LD')(D'^TL^T) \\
    & = (LD')(LD')^T \\
\end{align*}
Since $L$ is lower riangular and $D'$ is diagonal, $LD'$ is lower triangular.

Suppose $A = LL^T$ for some lower triangular $L$. Then:
\begin{align*}
    A^T & = (LL^T)^T \\
    & = LL^T = A
\end{align*}
Therefore, $A$ is symmetric. Further, for $x \neq 0$:
\begin{align*}
    x^T(LL^T)x & = (L^Tx)^T(L^Tx) \\
    & = ||L^Tx||^2 > 0
\end{align*}
Therefore $A$ is positive definite, and so it is symmetric positive definite.
\subproblem{(b)}
\begin{lstlisting}
function L = cholesky(A)
    L = A;
    n = size(A, 1);
    for k = 1:n
        for i = 1:k - 1
            s = 0;
            for j = 1:i - 1
                s = s + L(i, j) * L(k, j);
            end
            L(k, i) = (L(k, i) - s) / L(i, i);
        end
        v = 0;
        for j = 1:k - 1
            v = v + L(k, j) ^ 2;
        end
        L(k, k) = sqrt(L(k, k) - v);
        for i = k+1:n
            L(k, i) = 0;
        end
    end
end
\end{lstlisting}
\subproblem{(c)}
For $n = 3$, $A$ and $L$ are:
\begin{align*}
    A & = \begin{bmatrix}
        1.0000 & 0.5000 & 0.3333 \\
        0.5000 & 0.3333 & 0.2500 \\
        0.3333 & 0.2500 & 0.2000 \\
    \end{bmatrix} \\
    L & = \begin{bmatrix}
        1.0000 & 0.0000 & 0.0000 \\
        0.5000 & 0.2887 & 0.0000 \\
        0.3333 & 0.2887 & 0.0745 \\
    \end{bmatrix}
\end{align*}

For $n = 4$, $A$ and $L$ are:
\begin{align*}
    A & = \begin{bmatrix}
        1.0000 & 0.5000 & 0.3333 & 0.2500 \\
        0.5000 & 0.3333 & 0.2500 & 0.2000 \\
        0.3333 & 0.2500 & 0.2000 & 0.1667 \\
        0.2500 & 0.2000 & 0.1667 & 0.1429
    \end{bmatrix} \\
    L & = \begin{bmatrix}
        1.0000 & 0.0000 & 0.0000 & 0.0000 \\
        0.5000 & 0.2887 & 0.0000 & 0.0000 \\
        0.3333 & 0.2887 & 0.0745 & 0.0000 \\
        0.2500 & 0.2598 & 0.1118 & 0.0189 \\
    \end{bmatrix}
\end{align*}

For $n = 5$, $A$ and $L$ are:
\begin{align*}
    A & = \begin{bmatrix}
        1.0000 & 0.5000 & 0.3333 & 0.2500 & 0.2000 \\
        0.5000 & 0.3333 & 0.2500 & 0.2000 & 0.1667 \\
        0.3333 & 0.2500 & 0.2000 & 0.1667 & 0.1429 \\
        0.2500 & 0.2000 & 0.1667 & 0.1429 & 0.1250 \\
        0.2000 & 0.1667 & 0.1429 & 0.1250 & 0.1111 \\
    \end{bmatrix} \\
    L & = \begin{bmatrix}
        1.0000 & 0.0000 & 0.0000 & 0.0000 & 0.0000 \\
        0.5000 & 0.2887 & 0.0000 & 0.0000 & 0.0000 \\
        0.3333 & 0.2887 & 0.0745 & 0.0000 & 0.0000 \\
        0.2500 & 0.2598 & 0.1118 & 0.0189 & 0.0000 \\
        0.2000 & 0.2309 & 0.1278 & 0.0378 & 0.0048 \\
    \end{bmatrix}
\end{align*}

\clearpage

\problem{}
Let $||\cdot||$ be a vector norm on $\mathbb{R}^n$ and let $A \in \mathbb{R}^{n\times n}$ such that $\mathrm{rank}(A) = n$. Let $||x||_A = ||Ax||$. Then, for $x, y \in \mathbb{R}^n$:

Since $|| \cdot ||$ is a vector norm, $||x||_A = ||Ax|| > 0$ and $||\mathbf{0}|| = ||A \mathbf{0}|| = 0$. Further, for $\alpha \in \mathbb{R}$, $||\alpha x||_A = ||A(\alpha x)|| = |\alpha| \cdot ||Ax|| = |\alpha| \cdot ||x||_A$. Finally, $||A(x + y)|| = ||Ax + Ay|| \leq ||Ax|| + ||Ay||$, so $||x + y||_A \leq ||x||_A + ||y||_A$. Therefore $||x||_A$ is a vector norm on $\mathbb{R}^n$.

\clearpage

\problem{}
\subproblem{(a)}
Let $A$, $x$, $\hat{x}$, $b$, $\hat{b}$ such that $Ax = b$, $A\hat{x} = \hat{b}$, and $A$ is invertible. Then:
\begin{align*}
    A(x - \hat{x}) & = b - \hat{b} \\
    x - \hat{x} & = A^{-1}(b - \hat{b}) \\
    ||x - \hat{x}||_\infty & = ||A^{-1}(b - \hat{b})||_\infty \\
    & \leq ||A^{-1}||_\infty \cdot ||b - \hat{b}||_\infty
\end{align*}
Additionally, note that:
\begin{align*}
    ||Ax||_\infty & = ||b||_\infty \\
    ||A||_\infty \cdot ||x||_\infty & \geq \\
    \frac{1}{||x||_\infty} & \leq \frac{||A||_\infty}{||b||_\infty}
\end{align*}
Combining the inequalities:
\begin{align*}
    \frac{||x - \hat{x}||_\infty}{||x||_\infty} & \leq ||A||_\infty||A^{-1}||_\infty \cdot \frac{||b - \hat{b}||_\infty}{||b||_\infty} \\
    & \leq \kappa(A)\frac{||b - \hat{b}||_\infty}{||b||_\infty}
\end{align*}
\subproblem{(b)}
\begin{align*}
    A & = \begin{bmatrix}
        1 & 1.01 \\ 0.99 & 1
    \end{bmatrix} \\
    b & = \begin{bmatrix}
        2.01 \\ 1.99
    \end{bmatrix} \\
    \hat{b} & = \begin{bmatrix}
        2 \\ 2
    \end{bmatrix}
\end{align*}
The matrix $A$ is very close to singular; the inverse $A^{-1}$ is:
\begin{align*}
    A^{-1} & = \frac{1}{1 \cdot 1 - 1.01 \cdot 0.99}\begin{bmatrix}
        1 & -1.01 \\ -0.99 & 1
    \end{bmatrix} \\
    & = \begin{bmatrix}
        10000 & -10100 \\ -9900 & 10000
    \end{bmatrix}
\end{align*}
While $||A||_\infty = 2.01$, $||A^{-1}||_\infty = 20100$ and so $\kappa(A) = 40401$. This indicates that the relative error in $x$ can be up to $40401$ times the relative error in $b$. Even for a small perturbation of $0.01$ in $b$, this can result in a very large perturbation in $A$. Computing $x$ and $\bar{x}$:
\begin{align*}
    x & = \begin{bmatrix}
        1 \\ 1
    \end{bmatrix} \\
    \hat{x} & = \begin{bmatrix}
        -200 \\ 200
    \end{bmatrix}
\end{align*}
While the actual perturbation in $x$ is not as large as the upper bound, it is still hundreds of times larger than the perturbation in $b$.
\end{document}